\documentclass[11pt]{article}

% ========= 宏包 =========
\usepackage[a4paper,margin=1in]{geometry}
\usepackage{amsmath,amssymb,amsthm,mathtools}
\usepackage{hyperref}
\usepackage{enumitem}
\usepackage{ctex} % 中文支持（如你用 XeLaTeX/LuaLaTeX 编译）

% ========= 定理环境 =========
\newtheorem{definition}{定义}[section]
\newtheorem{theorem}[definition]{定理}
\newtheorem{proposition}[definition]{命题}
\newtheorem{lemma}[definition]{引理}
\newtheorem{example}[definition]{例}
\newtheorem{corollary}[definition]{推论}
\newtheorem{remark}[definition]{注}

% ========= 常用记号 =========
\newcommand{\R}{\mathbb{R}}
\newcommand{\D}{\mathcal{D}}
\newcommand{\X}{\mathfrak{X}}
\newcommand{\E}{\mathcal{E}} 
\newcommand{\supp}{\operatorname{supp}}

\title{积分曲线与流1}
\author{}
\date{}

\begin{document}
\maketitle
\tableofcontents

\section{积分曲线、流与流域}

\begin{definition}[积分曲线]
设 $M$ 为光滑流形，$V\in\X(M)$ 为光滑向量场。一条定义在开区间 $I\subset\R$ 上的光滑曲线
\[
\gamma:I\to M
\]
称为 $V$ 的\textbf{积分曲线}，若满足
\[
\gamma'(t)=V_{\gamma(t)}\quad(\forall t\in I).
\]
若它不能作为积分曲线延拓到更大的开区间，则称为\textbf{极大（maximal）积分曲线}。
\end{definition}

\begin{remark}[直观]
积分曲线描述\textbf{一个起点}在向量场“速度规则”下随时间演化的轨迹。
\end{remark}

\begin{definition}[流域与纤维时间集]
一个\textbf{流域}是开集 $\D\subseteq \R\times M$，使得对每个 $p\in M$，
\[
\D^{(p)}:=\{t\in\R:(t,p)\in\D\}
\]
是包含 $0$ 的开区间。
\end{definition}

\begin{definition}[（局部）流]
称映射 $\theta:\D\to M$ 为 $M$ 上的\textbf{流}，若 $\D$ 为流域且满足
\begin{align}
\theta(0,p)&=p,\quad \forall p\in M, \tag{9.6}\\
\theta(t,\theta(s,p))&=\theta(t+s,p)\quad \text{（当两边都有意义时）}. \tag{9.7}
\end{align}
若 $\D=\R\times M$，则称为\textbf{全局流}。
\end{definition}

\begin{remark}[直观]
流把\textbf{所有起点}的积分曲线统一打包成一个二元映射 $(t,p)\mapsto \theta(t,p)$，并满足“时间可加”的群律。
\end{remark}

\begin{definition}[时间切片 $M_t$]
给定流 $\theta:\D\to M$，对每个 $t\in\R$ 定义
\[
M_t:=\{p\in M:(t,p)\in\D\}. \tag{9.8}
\]
等价地，
\[
p\in M_t \iff t\in \D^{(p)} \iff (t,p)\in\D.
\]
\end{definition}

\begin{remark}[最关键的直观]
$M_t$ \textbf{不是}“时间 $t$ 时刻的点集”，而是“哪些\textbf{起点}允许演化时间 $t$”的集合（\textbf{定义域切片}）。
\end{remark}

\begin{definition}[无穷小生成元]
若 $\theta$ 是\textbf{光滑流}，其\textbf{无穷小生成元}定义为向量场
\[
V_p := \left.\frac{d}{dt}\right|_{t=0}\theta(t,p),
\]
其中 $\theta^{(p)}(t):=\theta(t,p)$。
\end{definition}

\begin{proposition}[命题 9.11（流产生积分曲线）]
若 $\theta:\D\to M$ 为光滑流，$V$ 为其无穷小生成元，则对每个 $p\in M$，
\[
t\mapsto \theta(t,p)
\]
是 $V$ 的积分曲线。
\end{proposition}

\begin{remark}[直观]
流是“整体运动”，生成元是“瞬时速度场”。
\end{remark}

\subsection{全局流的典型例子（例 9.8）}

\begin{example}[例 9.8(a)：平移流]
在 $M=\R^2$（坐标 $(x,y)$）上取 $V=\frac{\partial}{\partial x}$，则全局流为
\[
\tau_t(x,y)=(x+t,\;y).
\]
直观：$t>0$ 向右平移 $t$；$t<0$ 向左平移 $|t|$。
\end{example}

\begin{example}[例 9.8(b)：旋转流]
令
\[
W=x\frac{\partial}{\partial y}-y\frac{\partial}{\partial x},
\]
则全局流为
\[
\theta_t(x,y)=\bigl(x\cos t-y\sin t,\;x\sin t+y\cos t\bigr),
\]
直观：绕原点旋转角度 $t$（$t>0$ 逆时针）。
\end{example}

\subsection{并非所有向量场都有全局流（例 9.9--9.10 的直观）}

\begin{example}[例 9.9：流形“缺洞”导致有限时间无法延拓]
取 $M=\R^2\setminus\{0\}$，$V=\frac{\partial}{\partial x}$。从 $(-1,0)$ 出发的轨迹沿 $x$ 轴前进，会在有限时间碰到被删掉的点 $(0,0)$，因此不能作为 $M$ 内曲线继续延拓。
\end{example}

\begin{example}[例 9.10：有限时间“爆炸/逃逸到无穷远”]
取 $M=\R^2$，$W=x^2\frac{\partial}{\partial x}$。从 $(1,0)$ 出发，解在有限时间变为无界，导致不能延拓到所有时间。
\end{example}

\section{流的基本定理（定理 9.12）}

\begin{theorem}[定理 9.12：基本定理（极大流的存在唯一性与结构）]
设 $V$ 为 $M$ 上光滑向量场，则存在唯一的光滑\textbf{极大流} $\theta:\D\to M$，其生成元为 $V$，并且：
\begin{enumerate}[label=(\alph*)]
\item 对每个 $p\in M$，$\theta^{(p)}:\D^{(p)}\to M$ 是从 $p$ 出发的唯一极大积分曲线。
\item 若 $s\in\D^{(p)}$，则
\[
\D^{(\theta(s,p))}=\D^{(p)}-s.
\]
\item 对每个 $t\in\R$，$M_t$ 是开集，且
\[
\theta_t: M_t\to M_{-t}
\]
是微分同胚，逆为 $\theta_{-t}$。
\end{enumerate}
\end{theorem}

\begin{remark}[直观解释（不需要证明）]
\begin{itemize}
\item (a)：“固定起点 $p$”，$\theta(t,p)$ 就是那条积分曲线。
\item (b)：把起点沿流走到时间 $s$ 后，新起点能走的时间区间只是原区间整体平移。
\item (c)：$\theta_t$ 只推进了时间 $t$；落在 $M_{-t}$ 的意思是“到了像点后\textbf{允许再走回去} $-t$”，体现可逆性。注意 $M_t/M_{-t}$ 是\textbf{定义域切片}，不是“时刻层”，因此不存在“差了 $2t$ 的真实演化”这一说法。
\end{itemize}
\end{remark}

\section{流的自然性（命题 9.13 与推论 9.14）}

\begin{definition}[$F$-相关向量场]
设 $F:M\to N$ 光滑。若 $X\in\X(M)$、$Y\in\X(N)$ 满足
\[
dF_p(X_p)=Y_{F(p)}\quad(\forall p\in M),
\]
则称 $X$ 与 $Y$ 在 $F$ 下\textbf{相关（$F$-related）}。
\end{definition}

\begin{proposition}[命题 9.13：流的自然性]
设 $F:M\to N$ 光滑，$X\in\X(M)$、$Y\in\X(N)$，其流分别为 $\theta,\eta$。若 $X$ 与 $Y$ 在 $F$ 下相关，则对每个 $t\in\R$：
\[
F(M_t)\subseteq N_t,\qquad \eta_t\circ F=F\circ \theta_t\ \ \text{（在 }M_t\text{ 上）}.
\]
\end{proposition}

\begin{remark}[直观]
“先流再映射”等于“先映射再流”：动力学与映射 $F$ 相容（交换图）。
\end{remark}

\begin{corollary}[推论 9.14：微分同胚下流的不变性]
若 $F:M\to N$ 为微分同胚，$\theta$ 为 $X\in\X(M)$ 的流，则 $F_*X$ 的流为
\[
\eta_t=F\circ \theta_t\circ F^{-1},\qquad N_t=F(M_t).
\]
\end{corollary}

\section{完备向量场与判别准则（引理 9.15，定理 9.16--9.18，引理 9.19）}

\begin{definition}[完备（complete）]
向量场 $V$ 称为\textbf{完备}，若其流为\textbf{全局流}（即对所有 $t\in\R$、$p\in M$ 都有定义）；等价地，每条极大积分曲线定义域都是 $\R$。
\end{definition}

\begin{remark}[直观]
完备意味着“不会在有限时间因撞洞、逃逸到无穷远等原因而停止”。
\end{remark}

\begin{lemma}[引理 9.15：统一时间引理]
若存在 $\varepsilon>0$，使得对每个 $p\in M$，积分曲线 $\theta^{(p)}$ 的定义域都包含 $(-\varepsilon,\varepsilon)$，则 $V$ 完备。
\end{lemma}

\begin{remark}[直观]
每个点都能“至少走同样长的一小段时间”，就可以反复拼接，把解延拓到任意长时间，从而得到全局存在。
\end{remark}

\begin{theorem}[定理 9.16：紧支撑向量场完备]
光滑流形上每个\textbf{紧支撑}的光滑向量场都是完备的。
\end{theorem}

\begin{remark}[直观]
在支撑集外 $V=0$，轨道会变成常值曲线；结合紧致性与局部存在的“统一时间”，即可推出全局存在。
\end{remark}

\begin{corollary}[推论 9.17：紧流形上向量场完备]
在紧光滑流形上，每个光滑向量场都是完备的。
\end{corollary}

\begin{theorem}[定理 9.18：李群左不变向量场完备]
李群上的每个\textbf{左不变}向量场都是完备的。
\end{theorem}

\begin{remark}[直观]
左平移把单位元处的积分曲线搬运到任意点，因而得到对所有点统一的局部存在时间，进而推出完备。
\end{remark}

\begin{lemma}[引理 9.19：逃逸引理]
若 $\gamma:J\to M$ 为 $V$ 的极大积分曲线，且 $b=\sup J<\infty$，则对任意 $t_0\in J$，
\[
\gamma([t_0,b))\ \text{不包含在 }M\text{ 的任何紧子集中。}
\]
\end{lemma}

\begin{remark}[直观]
若解在有限时间 $b$ 处“走不下去”，轨道必然在逼近 $b$ 时逃离所有紧区域（跑向无穷远或撞到边界/缺口）。如果它一直待在某个紧集内，就能用统一时间把解延拓过 $b$，与极大性矛盾。
\end{remark}

\section{流出（Flowout）定理（定理 9.20）}

\begin{definition}[相切/处处不相切]
设 $S\subset M$ 为嵌入子流形，$V\in\X(M)$。
\begin{itemize}
\item $V$ 在 $p\in S$ 处\textbf{与 $S$ 相切}：$V_p\in T_pS$。
\item $V$ \textbf{处处不与 $S$ 相切}（nowhere tangent）：对所有 $p\in S$，$V_p\notin T_pS$。
\end{itemize}
若 $S=\{f=0\}$ 为超曲面（$df\neq 0$），则 $V_p\in T_pS \iff (Vf)(p)=0$。
\end{definition}

\begin{definition}[超曲面（hypersurface）]
若 $\dim M=n$，则 $S\subset M$ 称为\textbf{超曲面}，若 $S$ 是嵌入的 $(n-1)$ 维子流形（余维 $1$）。
等价地，局部可写为 $S=\{f=0\}$ 且在 $S$ 上 $df\neq 0$。
\end{definition}

\begin{theorem}[定理 9.20：流出定理]
设 $M$ 为光滑流形，$S\subseteq M$ 为嵌入的 $k$ 维子流形，$V\in\X(M)$ 处处不与 $S$ 相切。
令 $\theta:\D\to M$ 为 $V$ 的流，令
\[
\mathcal O=(\R\times S)\cap\D,\qquad \Phi=\theta|_{\mathcal O}.
\]
则：
\begin{enumerate}[label=(\alph*)]
\item $\Phi:\mathcal O\to M$ 是浸入（immersion）。
\item $\frac{\partial}{\partial t}\in\X(\mathcal O)$ 与 $V$ 在 $\Phi$ 下相关。
\item 存在光滑正函数 $\delta:S\to\R$，使得在
\[
\mathcal O_\delta=\{(t,p)\in\mathcal O:|t|<\delta(p)\} \tag{9.9}
\]
上 $\Phi$ 单射；因此 $\Phi(\mathcal O_\delta)$ 是包含 $S$ 的浸入子流形，且 $V$ 与之相切。
\item 若 $S$ 余维为 $1$（即超曲面），则 $\Phi|_{\mathcal O_\delta}$ 是到 $M$ 的某个开子流形上的微分同胚。
\end{enumerate}
\end{theorem}

\begin{remark}[直观]
把 $S$ 想成一张“横截面”。从 $S$ 上每点出发沿流走一小段时间，会“扫出”一条由轨道组成的管状邻域。参数 $t$ 控制沿轨道前进，$S$ 上的点（或其参数）用来给不同轨道编号。
\end{remark}

\section{正则点、奇点与平衡点（命题 9.21），以及正则点附近的规范形式（定理 9.22）}

\begin{definition}[奇点与正则点]
点 $p\in M$ 称为 $V$ 的\textbf{奇点}（singular），若 $V_p=0$；否则称为\textbf{正则点}（regular）。
\end{definition}

\begin{definition}[流的平衡点]
对流 $\theta:\D\to M$，点 $p$ 称为\textbf{平衡点}（equilibrium），若
\[
\theta(t,p)=p\quad \forall t\in\D^{(p)}.
\]
\end{definition}

\begin{proposition}[命题 9.21]
设 $V$ 的流为 $\theta:\D\to M$。
\begin{itemize}
\item 若 $p$ 是奇点，则 $\D^{(p)}=\R$ 且 $\theta^{(p)}(t)\equiv p$（常值曲线）。
\item 若 $p$ 是正则点，则 $\theta^{(p)}:\D^{(p)}\to M$ 是光滑浸入。
\end{itemize}
特别地，对于\textbf{光滑流}，其平衡点恰是无穷小生成元的奇点。
\end{proposition}

\begin{remark}[直观]
奇点处速度为零，所以点被“钉住”；正则点处速度不为零，轨道局部不会自我折返（浸入）。
\end{remark}

\begin{theorem}[定理 9.22：正则点附近的规范形式]
设 $V\in\X(M)$，$p\in M$ 为正则点。则存在 $p$ 的一组局部坐标 $(s^1,\dots,s^n)$ 使得
\[
V=\frac{\partial}{\partial s^1}.
\]
若给定超曲面 $S\ni p$ 且 $V_p\notin T_pS$，还可选坐标使得 $s^1$ 是 $S$ 的局部定义函数（局部上 $S=\{s^1=0\}$）。
\end{theorem}

\begin{remark}[直观]
在正则点附近，总能通过换坐标把向量场“拉直”为恒定的平移方向。于是局部动力学的复杂性只可能出现在平衡点/奇点附近。
\end{remark}

\subsection{规范形式的构造技巧（来自证明的“算法思路”）}

\begin{remark}[构造要点（无证明）]
\begin{enumerate}[label=\textbf{Step \arabic*:},leftmargin=2.5em]
\item 选取通过 $p$ 的超曲面 $S$，并保证 $V$ 在 $S$ 上处处不相切（必要时缩小邻域）。
\item 取 $S$ 的局部参数化 $X:\Omega\to S$（$\Omega\subset\R^{n-1}$）。
\item 定义“流出映射”
\[
\Psi(t,s)=\theta_t(X(s)).
\]
\item 对足够小的时间窗口与参数邻域，$\Psi$ 是到 $M$ 的某开集的微分同胚；于是用 $\Psi^{-1}$ 作为坐标图，其中 $s^1=t$。
\item 在该坐标里自动得到 $V=\partial/\partial s^1$，并且 $S$ 对应 $s^1=0$。
\end{enumerate}
\end{remark}

\subsection{例 9.23：旋转场在新坐标下变成 $\partial/\partial t$}

\begin{example}[例 9.23]
在 $M=\R^2$ 上令
\[
W=x\frac{\partial}{\partial y}-y\frac{\partial}{\partial x}.
\]
点 $p=(1,0)$ 为正则点。选取超曲面（曲线）
\[
S=\{(s,0): s \text{ 在 }1\text{ 附近}\}
\]
（$x$ 轴附近一段），参数化 $X(s)=(s,0)$。其流为旋转
\[
\theta_t(x,y)=(x\cos t-y\sin t,\;x\sin t+y\cos t).
\]
因此
\[
\Psi(t,s)=\theta_t(s,0)=(s\cos t,\ s\sin t).
\]
在 $(1,0)$ 邻域内可反解得到坐标
\[
(t,s)=\Psi^{-1}(x,y)=\left(\tan^{-1}\!\left(\frac{y}{x}\right),\ \sqrt{x^2+y^2}\right). \tag{9.13}
\]
在该坐标中 $W=\frac{\partial}{\partial t}$：它只改变角度 $t$，保持半径 $s$ 不变。
\end{example}

\begin{remark}关于“基向量变成什么”【补充直观】
旧坐标基为 $\partial_x,\partial_y$；新坐标基为 $\partial_t,\partial_s$。由于旋转运动保持 $s$（半径）不变，只改变 $t$（角度），故向量场在新坐标里只有 $\partial_t$ 分量，维数并未改变。
\end{remark}

\section{概念关系}

\begin{remark}
\begin{itemize}
\item 积分曲线：一个点的轨迹；流：所有点轨迹的打包，并满足群律。
\item 全局流 $\Leftrightarrow$ 向量场完备（所有时间都能演化）。
\item $M_t$ 是“允许走 $t$ 的起点集合”，不是“时刻 $t$ 的点集”。
\item 光滑情形：平衡点（流不动点） $\Leftrightarrow$ 奇点（$V_p=0$）。
\item 流出定理：从横截子流形 $S$ 出发，沿流扫出一个“轨道管状邻域”。
\item 规范形式定理：正则点附近总能换坐标把向量场拉直为 $\partial/\partial s^1$，复杂性集中在奇点附近。
\end{itemize}
\end{remark}

\section{李导数与可交换性}

\subsection{李导数（Lie Derivative）}

\begin{definition}[欧氏空间中的方向导数（回顾）]
在 $\R^n$ 上，令 $W$ 为光滑向量场，$v\in T_p\R^n\cong \R^n$。
定义 $W$ 在 $v$ 方向的方向导数为
\[
D_vW(p):=\left.\frac{d}{dt}\right|_{t=0}W(p+tv)
=\lim_{t\to 0}\frac{W(p+tv)-W(p)}{t}.
\tag{9.15}
\]
等价地（按分量逐一求导）：
若 $W=W^i\frac{\partial}{\partial x^i}$，则
\[
D_vW(p)=\bigl(D_vW^i(p)\bigr)\frac{\partial}{\partial x^i}\Big|_p.
\]
\end{definition}

\begin{remark}[为什么在流形上不能直接照搬 (9.15)]
在一般流形上，$W_{\gamma(t)}\in T_{\gamma(t)}M$ 与 $W_{\gamma(0)}\in T_{\gamma(0)}M$
属于不同切空间，不能直接相减；因此需要用流把向量“搬回同一切空间”再作差商。
\end{remark}

\begin{definition}[李导数：向量场对向量场]
设 $M$ 为光滑流形，$V,W\in\X(M)$，令 $\theta$ 为 $V$ 的流。
定义 $W$ 关于 $V$ 的\textbf{李导数} $L_VW$ 为如下（一般仅先得到粗糙向量场）：
对每个 $p\in M$，若导数存在，令
\[
(L_VW)_p
:=\left.\frac{d}{dt}\right|_{t=0} \, d(\theta_{-t})_{\theta_t(p)}\bigl(W_{\theta_t(p)}\bigr)
=\lim_{t\to 0}\frac{d(\theta_{-t})_{\theta_t(p)}(W_{\theta_t(p)})-W_p}{t}.
\tag{9.16}
\]
\end{definition}

\begin{remark}[几何直观]
先沿 $V$ 的流把点 $p$ 推到 $\theta_t(p)$，在那里取向量 $W_{\theta_t(p)}$；
再用反向流 $\theta_{-t}$ 的微分把该向量推回到 $T_pM$；
最后在同一切空间 $T_pM$ 中对 $t$ 求导。
\end{remark}

\begin{lemma}[引理 9.36：李导数存在且光滑]
设 $M$ 为带边界或无边界的光滑流形，$V,W\in\X(M)$。
若 $\partial M\neq\varnothing$，另假设 $V$ 与 $\partial M$ 相切。
则对每个 $p\in M$，$(L_VW)_p$ 存在，且 $L_VW$ 是光滑向量场。
\end{lemma}

\begin{theorem}[定理 9.38：李导数与李括号一致]
对任意光滑向量场 $V,W\in\X(M)$，
\[
L_VW=[V,W].
\]
\end{theorem}

\subsection{李导数（李括号）的性质}

\begin{corollary}[推论 9.39：李导数的基本恒等式]
设 $M$ 为光滑流形（可有边界），$V,W,X\in\X(M)$，则：
\begin{enumerate}[label=(\alph*)]
\item 反对称性：
\[
L_VW=-L_WV.
\]
\item 对李括号的导子性质：
\[
L_V[W,X]=[L_VW,X]+[W,L_VX].
\]
\item 交换子恒等式（导子交换）：
\[
L_{[V,W]}X=L_VL_WX-L_WL_VX.
\]
\item 乘积法则（函数与向量场）：
若 $g\in C^\infty(M)$，则
\[
L_V(gW)=(Vg)\,W+g\,L_VW.
\]
\item 与微分同胚的自然性：
若 $F:M\to N$ 是微分同胚，则
\[
F_*(L_VX)=L_{F_*V}(F_*X).
\]
\end{enumerate}
\end{corollary}

\subsection{在任意时刻的“搬回”导数公式}

\begin{proposition}[命题 9.41：任意时刻的导数公式]
设 $M$ 为光滑流形（可有边界），$V,W\in\X(M)$；
若 $\partial M\neq\varnothing$，假设 $V$ 与 $\partial M$ 相切。
令 $\theta$ 为 $V$ 的流。对任意 $(t_0,p)$ 在流的定义域内，有
\[
\left.\frac{d}{dt}\right|_{t=t_0}
d(\theta_{-t})_{\theta_t(p)}\bigl(W_{\theta_t(p)}\bigr)
=
d(\theta_{-t_0})_{\theta_{t_0}(p)}
\Bigl((L_VW)_{\theta_{t_0}(p)}\Bigr).
\tag{9.17}
\]
\end{proposition}

\subsection{可交换向量场与流的不变性}

\begin{definition}[可交换向量场]
设 $M$ 为光滑流形，$V,W\in\X(M)$。称 $V$ 与 $W$ \textbf{可交换（commute）}，
若对任意 $f\in C^\infty(M)$，
\[
VWf=WVf,
\]
等价地，
\[
[V,W]\equiv 0.
\]
\end{definition}

\begin{definition}[向量场在流下不变（invariant under a flow）]
设 $\theta$ 为 $M$ 上的光滑流。称向量场 $W$ 在 $\theta$ 下\textbf{不变}，
若对每个 $t$，$W$ 与自身在 $\theta_t$ 下相关；等价地，
对所有定义域内的 $(t,p)$，
\[
d(\theta_t)_p(W_p)=W_{\theta_t(p)}.
\]
\end{definition}

\begin{theorem}[定理 9.42：可交换 $\Longleftrightarrow$ 互相不变]
设 $V,W$ 是光滑流形 $M$ 上的光滑向量场。则以下命题等价：
\begin{enumerate}[label=(\alph*)]
\item $V$ 与 $W$ 可交换。
\item $W$ 在 $V$ 的流下不变。
\item $V$ 在 $W$ 的流下不变。
\end{enumerate}
\end{theorem}

\begin{corollary}[推论 9.43：每个向量场在自身流下不变]
对任意 $V\in\X(M)$，$V$ 在其自身的流下不变（等价地 $[V,V]\equiv 0$）。
\end{corollary}

\subsection{流的可交换性：局部定义域版本}

\begin{definition}[流可交换（非完备情形的定义）]
设 $\theta$ 与 $\psi$ 是 $M$ 上两个（可能仅局部定义的）流。
称 $\theta$ 与 $\psi$ \textbf{可交换}，若对每个 $p\in M$ 满足：
若存在包含 $0$ 的开区间 $J,K$ 使得
\[
\theta_t\circ\psi_s(p)\ \text{或}\ \psi_s\circ\theta_t(p)
\]
其中之一对所有 $(s,t)\in J\times K$ 都有定义，
则二者在 $J\times K$ 上都定义且满足
\[
\theta_t(\psi_s(p))=\psi_s(\theta_t(p)),\qquad \forall (s,t)\in J\times K.
\]
若两者为全局流，则该定义等价于
$\theta_t\circ\psi_s=\psi_s\circ\theta_t$ 对所有 $s,t\in\R$ 成立。
\end{definition}

\begin{theorem}[定理 9.44：向量场可交换 $\Longleftrightarrow$ 流可交换]
设 $V,W\in\X(M)$ 为光滑向量场，其流分别为 $\theta,\psi$。
则
\[
[V,W]\equiv 0
\quad\Longleftrightarrow\quad
\theta\ \text{与}\ \psi\ \text{可交换（按上一定义）}.
\]
\end{theorem}

\section{可交换标架与可交换向量场的规范形}

\subsection{可交换标架}

\begin{definition}[局部标架与可交换标架]
设 $M$ 为光滑 $n$ 维流形。$U\subseteq M$ 上的一个\textbf{局部标架}是
$n$ 元组向量场 $(E_i)$，使得对每个 $p\in U$，$\{E_i|_p\}$ 构成 $T_pM$ 的一组基。
若该标架满足
\[
[E_i,E_j]=0,\qquad \forall i,j,
\]
则称为\textbf{可交换标架}（亦称 holonomic frame）。
\end{definition}

\begin{remark}[坐标标架是可交换标架]
任意坐标图 $(U,(x^i))$ 的坐标标架 $\bigl(\partial/\partial x^i\bigr)$ 满足
$[\partial/\partial x^i,\partial/\partial x^j]=0$，因此必为可交换标架。
\end{remark}

\subsection{可交换向量场的标准形（Frobenius 的局部版本）}

\begin{theorem}[定理 9.46：可交换向量场的规范坐标（Canonical Form）]
设 $M$ 为光滑 $n$ 维流形，
$(V_1,\dots,V_k)$ 是开集 $W\subseteq M$ 上一组线性无关的
\textbf{两两可交换}的光滑向量场。
则对每个 $p\in W$，存在以 $p$ 为中心的光滑坐标图 $(U,(s^i))$ 使得
\[
V_i=\frac{\partial}{\partial s^i},\qquad i=1,\dots,k.
\]
进一步，若 $S\subseteq W$ 是嵌入的余维 $k$ 子流形，且 $p\in S$ 并满足
\[
T_pS\ \text{与}\ \mathrm{span}(V_1|_p,\dots,V_k|_p)\ \text{互补},
\]
则可选坐标使得
\[
S\cap U=\{s^1=\cdots=s^k=0\}.
\]
\end{theorem}

\section{时间依赖向量场与时间依赖流}

\subsection{时间依赖向量场}

\begin{definition}[时间依赖向量场]
设 $M$ 为光滑流形，$J\subseteq\R$ 为区间。
称连续映射
\[
V:J\times M\to TM
\]
为 $M$ 上的\textbf{时间依赖向量场}，若对每个 $(t,p)\in J\times M$，
\[
V(t,p)\in T_pM.
\]
对每个 $t\in J$，定义 $V_t\in\X(M)$ 为
\[
V_t(p):=V(t,p).
\]
\end{definition}

\begin{definition}[时间依赖向量场的积分曲线]
若 $V$ 为时间依赖向量场，则称可微曲线 $\gamma:J_0\to M$
（其中 $J_0\subseteq J$ 为区间）为 $V$ 的\textbf{积分曲线}，若满足
\[
\gamma'(t)=V(t,\gamma(t)),\qquad \forall t\in J_0.
\]
\end{definition}

\subsection{时间依赖流的基本定理}

\begin{theorem}[定理 9.48：时间依赖流的基本定理]
设 $M$ 为光滑流形，$J\subseteq\R$ 为开区间，
且 $V:J\times M\to TM$ 为光滑时间依赖向量场。
则存在开集 $\E\subseteq J\times J\times M$ 与光滑映射
\[
\psi:\E\to M,
\]
称为 $V$ 的\textbf{时间依赖流}，并满足：
\begin{enumerate}[label=(\alph*)]
\item 对每个 $t_0\in J$ 与 $p\in M$，
\[
\E^{(t_0,p)}:=\{t\in J:(t,t_0,p)\in\E\}
\]
是包含 $t_0$ 的开区间；并且曲线
\[
\psi^{(t_0,p)}:\E^{(t_0,p)}\to M,\qquad
\psi^{(t_0,p)}(t)=\psi(t,t_0,p),
\]
是满足初值条件 $\psi^{(t_0,p)}(t_0)=p$ 的唯一极大积分曲线。
\item 若 $t_1\in\E^{(t_0,p)}$ 且 $q=\psi^{(t_0,p)}(t_1)$，则
\[
\E^{(t_1,q)}=\E^{(t_0,p)},\qquad \psi^{(t_1,q)}=\psi^{(t_0,p)}.
\]
\item 对每个 $(t_1,t_0)\in J\times J$，集合
\[
M_{t_1,t_0}:=\{p\in M:(t_1,t_0,p)\in\E\}
\]
在 $M$ 中是开集，且映射
\[
\psi_{t_1,t_0}:M_{t_1,t_0}\to M,\qquad
\psi_{t_1,t_0}(p)=\psi(t_1,t_0,p),
\]
是从 $M_{t_1,t_0}$ 到 $M_{t_0,t_1}$ 的微分同胚，其逆为 $\psi_{t_0,t_1}$。
\item 若 $p\in M_{t_1,t_0}$ 且 $\psi_{t_1,t_0}(p)\in M_{t_2,t_1}$，
则 $p\in M_{t_2,t_0}$ 且满足复合律
\[
\psi_{t_2,t_1}\circ\psi_{t_1,t_0}(p)=\psi_{t_2,t_0}(p).
\tag{9.18}
\]
\end{enumerate}
\end{theorem}

\end{document}
